% Abstract for EpiContext paper (v7 - Breakthrough Results)
Large language model (LLM) agents face a fundamental tension in long-horizon tasks: maintaining comprehensive context history consumes prohibitive token budgets, while aggressive compression discards critical information. We introduce \EpiContext, a framework that reformulates agent context management through the lens of biological epigenetics. In \EpiContext, the agent's complete interaction history constitutes an immutable ``genome''---a structured context graph---while each request payload represents a dynamic ``epigenetic expression'' regulated by three operators: \textit{methylation} (silencing noisy memories), \textit{acetylation} (activating relevant tools), and \textit{fitness feedback} (adaptive threshold tuning). We implement \EpiContext{} as a custom agent within the Harbor evaluation framework and conduct controlled experiments comparing six context strategies across five containerized tasks. Our initial implementation (v1) underperformed simple baselines due to metadata overhead and conservative context retention. Through systematic root cause analysis and engineering improvements---compact activation encoding, adaptive strategy switching, and aggressive filtering---our improved Adaptive \EpiContext{} (v2) achieves 90\% token reduction compared to Full-Context ($p = 0.022$) and 64\% fewer turns ($p < 0.001$), establishing a new state-of-the-art in context efficiency among the compared strategies. On the most complex task, Adaptive \EpiContext{} reduces token consumption by 96\% (1,801 vs.\ 43,600) while completing the task in 5 turns instead of 20. Our work demonstrates that principled engineering of epigenetic context regulation can dramatically improve agent efficiency, and establishes a reproducible experimental framework for future investigation.